\section{Related Works}
\label{sec:related_work}


% \section{Related Works}

% \begin{itemize}
%     \item Comparing different model families
%     \begin{itemize}
%         \item OV: GroundingDINO, SAM
%         \item MLLMs
%         \item Visual Encoders
%     \end{itemize}
%     \item Text Conditioned Visual Feats
%     \begin{itemize}
%         \item FLAIR
%         \item SAM3 --> MaskInversion
%         \item ELIP, TIE
%     \end{itemize}
% \end{itemize}




\paragraph{Visual Representation Families.}
\JR{in the rest of the paper we use bold text instead of italics for paragraph headings (i.e. subsubsections)}
\YA{technically we should also look at CoCoOp}
Unimodal encoders (DINOv2~\cite{oquab2024dinov2}, MAE~\cite{he2022mae}) learn rich visual representations via self-supervised objectives but produce entirely query-agnostic features.
Cross-modal encoders (CLIP~\cite{radford2021_clip}, SigLIP~\cite{zhai2023_siglip}, CoCoOp~\cite{zhou2022_cocoop}) align image--text embeddings via contrastive or prompt-based pretraining but limit multimodal interaction to post-hoc fusion.
MLLMs achieve moderate steerability by routing visual tokens through a language backbone at substantial computational cost~\cite{chen2025vlm_benchmark}, while open-vocabulary models (GroundingDINO~\cite{liu2023groundingdino}, SAM3~\cite{carion2025_sam3}) are highly steerable but produce localization-specific features that transfer poorly.
We compare these families quantitatively in \cref{tab:method_comparison} and \cref{fig:teaser}.

%%% === Old Visual Representation Families (commented out) ===
% {\color{gray} Modern vision systems extract features through architecturally distinct paradigms.}
% \textit{Unimodal encoders} such as DINOv2~\cite{oquab2024dinov2} and MAE~\cite{he2022mae} learn rich visual representations via self-supervised objectives but produce features that are entirely query-agnostic.
% \textit{Cross-modal encoders} like CLIP~\cite{radford2021_clip} and SigLIP~\cite{zhai2023_siglip} align image and text embeddings through contrastive pretraining, yet their dual-encoder design limits multimodal interaction to post-hoc fusion rather than early feature-level conditioning. 
% \textit{Multimodal large language models} (MLLMs) process interleaved image--text sequences and can produce text-influenced visual representations, but they route visual tokens through a large language backbone, incurring substantial computational overhead for moderate steerability.
% \textit{Open-vocabulary localization models} such as GroundingDINO~\cite{liu2023groundingdino} and SAM3~\cite{carion2025_sam3} fuse text deeply into their encoder through cross-modal attention, yielding highly steerable intermediate features; however, these representations are optimized for localization and transfer poorly to general vision tasks.
% Our work bridges this gap: we condition frozen ViT features on text within the transformer layers themselves, achieving high steerability while preserving the representation quality of the base encoder.
%%% === End old Visual Representation Families ===


\paragraph{Text-Conditioned Visual Features.}
\YA{we should say that this family basically doesn't exist. and highly differentiate/dismiss FLAIR }
% TODO: expand with FLAIR, SAM3/MaskInversion, ELIP, TIE
The goal of producing visual features that are both text-steerable and high-quality remains largely unaddressed.
The closest attempt, FLAIR~\cite{xiao2025flair}, applies text-conditioned attention pooling over frozen CLIP features.
Because this fusion occurs outside the encoder, at the pooling stage, text cannot influence the encoding process itself: FLAIR achieves only moderate steerability and weaker text-dependence (\cref{subsec:cond_ret_exp}), and its features fall short of unimodal encoders on standard vision benchmarks (\cref{fig:teaser}).
Moreover, FLAIR is restricted to CLIP-style backbones and requires two orders of magnitude more training data.
These limitations reflect a deeper asymmetry: high-quality unimodal data---both vision-only and language-only---vastly exceeds paired multimodal data, which is why unimodal encoders learn richer representations in the first place.
This suggests that rather than training joint vision-language encoders from scratch on limited paired data, one should inject language conditioning into strong unimodal vision encoders, preserving the representational depth learned from each modality independently.
The challenge then becomes one of \textit{representation alignment}: how to merge two unimodal experts into an early-fused encoder without degrading either modality's representations.
\model{} addresses this by conditioning a frozen ViT on text within the transformer layers themselves; text serves not as a training signal but as a steering mechanism that preserves the base encoder's quality while extending generalization to domains where labeled data is scarce.

%%% === Old Text-Conditioned Visual Features (commented out) ===
% FLAIR~\cite{xiao2025flair} applies text-conditioned attention pooling over frozen CLIP features to produce text-informed global representations. Unlike \model{}, it operates at the pooling stage rather than within the encoder itself.
% Failure of FLAIR -->
%%% === End old Text-Conditioned Visual Features ===

% SAM3 --> MaskInversion
% ELIP, TIE

